\chapter{Pendahuluan}

\section{Latar Belakang}
Tuliskanlah latar belakang dari pelaksanaan kerja praktik (KP) di perusahaan dan substansi yang digeluti berkaitan dengan tujuan, misi, visi atau fungsi perusahaan. Jangan menuliskan tentang pelaksanaan mata kuliah Kerja Praktek seperti yang ditulis di dalam kurikulum.

Untuk menuliskan paragraf berikutnya berikan jarak satu Enter seperti contoh pada template ini. Sebuah paragraf hendaknya mengandung satu pokok pikiran dan beberapa kalimat gagasan pendukung.

\section{Tempat dan Waktu Pelaksanaan}
Tuliskan tempat atau lokasi dan waktu pelaksanaan kegiatan kerja praktik.

\section{Ruang Lingkup Pekerjaan}
Tuliskan ruang lingkup pekerjaan di tempat KP, bukan lingkup KP seperti dituliskan pada kurikulum. Deskripsi pekerjaan secara lengkap dapat dijelaskan di Bab III.

\section{Sistematika Penulisan}
Tuliskanlah sistematika penulisan laporan hasil kerja praktik ini yang dapat mencakup hal-hal sebagai berikut:

\begin{enumerate}[label=\arabic*)]
\item Bab I Pendahuluan\\
Bab ini \ldots
\item Bab II Organisasi atau Lingkungan Kerja Praktik\\
Bab ini berisi \ldots
\item Bab III Judul Pekerjaan\\
Bab ini berisi \ldots
\item Bab IV Rangkuman\\
Bab ini berisi \ldots
\end{enumerate}
